\chapter{Memory Kernel}

\section{The Foundation Layer}

The Memory Kernel is the bedrock of the entire Connector platform. Every other
layer — execution engines, protocol stacks, developer tooling — sits on top of
it. Its job is deceptively simple: \textbf{ensure that all agent state is
persistent, immutable, addressable by content, and cryptographically
verifiable}. Nothing happens in Connector that is not mediated by the kernel.

The best analogy is an operating system kernel combined with a content-addressed
filesystem. An OS kernel mediates every system call a process makes; Connector's
Memory Kernel mediates every operation an agent performs. An OS filesystem
tracks file paths and permissions; the Memory Kernel tracks content hashes and
cryptographic lineage.

\begin{insightbox}{Why content addressing?}
Traditional databases identify records by primary key — an arbitrary number that
tells you nothing about the data. Content addressing identifies records by their
\emph{content hash}. If you have the CID, you can verify the content is
unchanged without trusting any server. Identical content produces identical CIDs,
giving automatic deduplication. This is how Connector provides tamper-evident
audit trails without a separate audit database.
\end{insightbox}

\section{The MemPacket}

A \textbf{MemPacket} is the atomic unit of all state in Connector. Every piece
of information — an agent's reasoning output, a tool call result, a compliance
record, a session checkpoint — is stored as a MemPacket. There is no other
storage primitive.

\begin{figure}[H]
\centering
\begin{tikzpicture}[
  every node/.style={font=\small},
  field/.style={rectangle, draw=MidnightBlue!50, fill=MidnightBlue!6,
                text width=3.8cm, align=center, minimum height=0.75cm,
                inner sep=4pt},
  label/.style={font=\footnotesize\itshape, text=gray, align=left},
  node distance=0.3cm
]

%% Outer packet box
\node[rectangle, rounded corners=6pt, draw=MidnightBlue, fill=MidnightBlue!4,
      minimum width=10cm, minimum height=6.8cm, label={[font=\bfseries\small,
      text=MidnightBlue]above:MemPacket}] (packet) at (0,0) {};

%% Fields
\node[field] (cid) at (0, 2.6)     {\texttt{cid}\\\footnotesize cls1-sha256-<64hex>};
\node[field] (content) at (0, 1.7) {\texttt{content}\\\footnotesize Typed union (8 variants)};
\node[field] (meta) at (0, 0.8)    {\texttt{metadata}\\\footnotesize Timestamps, tags, lineage};
\node[field] (sig) at (0,-0.1)     {\texttt{signature}\\\footnotesize Ed25519 (optional)};
\node[field] (enc) at (0,-1.0)     {\texttt{encryption}\\\footnotesize Envelope (optional)};

%% Content variants panel
\node[rectangle, rounded corners=4pt, draw=CExec!60, fill=CExec!6,
      text width=3cm, align=center, minimum height=3.8cm,
      font=\footnotesize, inner sep=6pt]
  (variants) at (4.4, 1.2) {
    \textbf{Content Types}\\[4pt]
    Text\\Structured\\Binary\\Perception\\Meaning\\Thought\\Action\\Evidence
  };

\draw[-Stealth, dashed, CExec!60] (content.east) -- (variants.west);

%% Lineage annotation
\node[label, right=0.2cm of meta.east, text width=2.8cm] (linann)
  {parent\_cids: [\ldots]\\session\_id\\agent\_did\\tags};

\end{tikzpicture}
\caption{MemPacket anatomy — the atomic state unit of Connector}
\label{fig:mempacket}
\end{figure}

\subsection{Content Types}

The \texttt{content} field is a typed union with eight variants, covering the
full range of information an agent might produce:

\begin{table}[H]
\centering
\renewcommand{\arraystretch}{1.35}
\begin{tabularx}{\textwidth}{@{}llX@{}}
\toprule
\textbf{Type} & \textbf{Use case} & \textbf{Typical producer} \\
\midrule
\texttt{Text}        & Raw text, prompts, responses      & LLM gateway \\
\texttt{Structured}  & JSON data, query results          & Tool calls, API responses \\
\texttt{Binary}      & Files, images, embeddings         & Asset pipeline \\
\texttt{Perception}  & Sensor input, multimodal data     & Cognitive substrate Layer 1 \\
\texttt{Meaning}     & Semantic extraction result        & Cognitive substrate Layer 2 \\
\texttt{Thought}     & Deliberation record               & Cognitive substrate Layers 3--10 \\
\texttt{Action}      & Tool call, memory write record    & Execution engine \\
\texttt{Evidence}    & Audit chain, compliance record    & Kernel, SOE layer \\
\bottomrule
\end{tabularx}
\caption{MemPacket content type variants and their producers}
\end{table}

\section{Content Addressing and CIDs}

Every MemPacket is identified by its \textbf{Content Identifier (CID)}, computed
deterministically from the packet's content:

\[
\texttt{CID} = \texttt{cls1-sha256-} \;\|\; \text{SHA256}\!\left(\text{canonical\_json}(\textit{content})\right)
\]

The prefix \texttt{cls1-sha256-} identifies the Connector namespace and hash
algorithm, making CIDs self-describing and forward-compatible.

\begin{figure}[H]
\centering
\begin{tikzpicture}[
  box/.style={rectangle, rounded corners=3pt, draw=#1!70, fill=#1!10,
              align=center, font=\footnotesize\bfseries, inner sep=6pt,
              minimum height=0.8cm},
  node distance=0.9cm and 1.1cm
]

\node[box=CExec, text width=3cm] (content) {Agent Output\\(any content)};
\node[box=CProto, text width=3.2cm, right=of content] (json)
  {Canonical JSON\\serialisation};
\node[box=CKernel, text width=3.2cm, right=of json] (hash)
  {SHA-256\\hash};
\node[box=CKernel, text width=3.4cm, right=of hash] (cid)
  {\texttt{cls1-sha256-}\\$\langle$64 hex chars$\rangle$};

\draw[arr=gray] (content) -- (json);
\draw[arr=gray] (json) -- (hash);
\draw[arr=gray] (hash) -- (cid);

\node[below=0.5cm of json, font=\footnotesize\itshape, text=gray]
  {deterministic, field-sorted};
\node[below=0.5cm of hash, font=\footnotesize\itshape, text=gray]
  {32 bytes = 256 bits};
\node[below=0.5cm of cid, font=\footnotesize\itshape, text=gray]
  {self-describing prefix};

\end{tikzpicture}
\caption{CID computation pipeline — content → canonical JSON → SHA-256 → CID}
\label{fig:cid-flow}
\end{figure}

Four properties flow directly from this design:

\begin{enumerate}
  \item \textbf{Deduplication} — Two packets with identical content share a CID.
        The storage layer never writes the same data twice.
  \item \textbf{Integrity} — Any modification to content changes the CID.
        Corruption is detected without a separate checksum database.
  \item \textbf{Independent verification} — Anyone with the content can recompute
        the CID and verify it matches, without trusting the server.
  \item \textbf{Immutability} — MemPackets are write-once. Updating data means
        writing a new packet that references the old CID as its parent, creating
        an auditable lineage chain.
\end{enumerate}

\section{Namespace Hierarchy}

The kernel organises memory into nine namespace types, each with distinct
security levels and access semantics, analogous to Linux filesystem mounts:

\begin{table}[H]
\centering
\renewcommand{\arraystretch}{1.4}
\begin{tabular}{@{}cllll@{}}
\toprule
\textbf{Prefix} & \textbf{Type} & \textbf{Linux Analogy} & \textbf{Security} & \textbf{Purpose} \\
\midrule
\texttt{/m/} & Memory     & \texttt{/home} & Standard & Private agent working memory \\
\texttt{/k/} & Knowledge  & \texttt{/usr}  & Protected & Shared, validated knowledge bases \\
\texttt{/v/} & Asset      & \texttt{/var}  & Tool I/O & Raw input files, staging area \\
\texttt{/x/} & App        & \texttt{/opt}  & Tool I/O & Executable apps via MCP \\
\texttt{/c/} & Core       & \texttt{/bin}  & Control  & Native tools via CNP \\
\texttt{/a/} & Agent      & \texttt{/proc} & Control  & Agent control plane \\
\texttt{/t/} & Tool       & \texttt{/tmp}  & Tool I/O & Ephemeral tool I/O \\
\texttt{/s/} & System     & \texttt{/sys}  & Kernel   & Kernel-reserved, read-only \\
\texttt{/p/} & Public     & \texttt{/srv}  & Public   & Announcements, shared outputs \\
\bottomrule
\end{tabular}
\caption{Memory namespace types and their security boundaries}
\end{table}

\section{Syscall Interface}

The kernel exposes over 60 syscalls organised into eight categories.
Rather than raw code, what matters is the contract each category provides:

\begin{figure}[H]
\centering
\begin{tikzpicture}[
  cat/.style={rectangle, rounded corners=4pt, draw=MidnightBlue!50,
              fill=MidnightBlue!7, text width=4.8cm, align=left,
              minimum height=1.1cm, font=\footnotesize, inner sep=6pt},
  node distance=0.45cm and 0.7cm
]
\node[cat] (mem)  {\textbf{Memory R/W}\\\texttt{Write, Read, Query, Search}};
\node[cat, right=of mem] (agent) {\textbf{Agent Lifecycle}\\\texttt{Register, Start, Pause, Terminate}};
\node[cat, below=of mem] (sess) {\textbf{Session Mgmt}\\\texttt{Create, Close, Checkpoint, Restore}};
\node[cat, right=of sess] (cap)  {\textbf{Capabilities}\\\texttt{Grant, Revoke, Attenuate, Verify}};
\node[cat, below=of sess] (pred) {\textbf{Predicate Eval}\\\texttt{Check, Enforce, Audit}};
\node[cat, right=of pred] (lin)  {\textbf{Lineage}\\\texttt{GetLineage, GetChildren, Trace}};
\node[cat, below=of pred] (cry)  {\textbf{Cryptography}\\\texttt{Sign, Verify, Seal, Open}};
\node[cat, right=of cry] (tool)  {\textbf{Tool Bridge}\\\texttt{Invoke, List, Grant, Revoke}};
\end{tikzpicture}
\caption{Eight syscall categories exposed by the Memory Kernel}
\label{fig:syscalls}
\end{figure}

\section{Agent Lifecycle}

Every agent in Connector passes through a well-defined state machine. The kernel
enforces these transitions — no agent can execute outside an \texttt{Active}
state, and no agent exits \texttt{Active} without a checkpoint.

\begin{figure}[H]
\centering
\begin{tikzpicture}[
  state/.style={circle, draw=MidnightBlue!70, fill=MidnightBlue!10,
                minimum size=1.8cm, align=center, font=\small\bfseries},
  term/.style={circle, draw=CDanger!70, fill=CDanger!10,
               minimum size=1.8cm, align=center, font=\small\bfseries},
  init/.style={circle, draw=CGreen!70, fill=CGreen!10,
               minimum size=1.8cm, align=center, font=\small\bfseries},
  node distance=2.4cm
]

\node[init]  (reg) {Registered};
\node[state, right=of reg]  (act) {Active};
\node[state, below=of act]  (pau) {Paused};
\node[term,  right=of act]  (ter) {Terminated};

\draw[arr=CGreen] (reg) -- node[above,font=\footnotesize]{\texttt{start}} (act);
\draw[arr=MidnightBlue] (act) -- node[right,font=\footnotesize]{\texttt{pause}} (pau);
\draw[arr=MidnightBlue] (pau.north) to[bend right=20]
  node[left,font=\footnotesize]{\texttt{resume}} (act.south);
\draw[arr=CDanger] (act) -- node[above,font=\footnotesize]{\texttt{terminate}} (ter);
\draw[arr=CDanger] (pau) to[bend right=15]
  node[below right,font=\footnotesize]{\texttt{terminate}} (ter);

%% Checkpoint annotations
\node[below=0.3cm of pau, font=\footnotesize\itshape, text=gray]
  {kernel writes checkpoint};
\node[above=0.3cm of act, font=\footnotesize\itshape, text=gray]
  {predicates enforced on every op};

\end{tikzpicture}
\caption{Agent lifecycle state machine enforced by the Memory Kernel}
\label{fig:agent-lifecycle}
\end{figure}

A key constraint: the transition from \texttt{Active} to any other state
\emph{always} triggers a checkpoint write. This means that even an unexpected
crash (which skips the \texttt{terminate} transition) leaves the kernel able to
restore the agent to its last checkpoint on the next boot.

\section{Session Management}

An \textbf{Agent} is a long-lived identity; a \textbf{Session} is a scoped
execution context within that agent. Think of the agent as a Unix process and
the session as a single system call batch — isolated, budgeted, and auditable.

\begin{conceptbox}{Session Contract}
A session is created with three inputs, and produces one output:
\begin{itemize}[noitemsep]
  \item \textbf{Budget} — Maximum tokens, memory bytes, tool calls, wall time
  \item \textbf{Predicates} — Governance rules that must hold throughout execution
  \item \textbf{Agent DID} — The identity executing within this session
\end{itemize}
Output: a signed \textbf{ExecutionReceipt} containing the full operation log,
final state CID, and budget consumption.
\end{conceptbox}

The session lifecycle is: \textbf{Create → Execute → Checkpoint* → Close → Receipt}.
Checkpointing can occur multiple times during execution. If the node crashes
mid-session, the session can be restored from the last checkpoint.

\section{The 11-Layer Predicate Chain}

Before any kernel operation executes, the request passes through eleven
enforcement layers in sequence. If any layer rejects the operation, execution
stops immediately and an audit entry is written.

\begin{figure}[H]
\centering
\begin{tikzpicture}[
  layer/.style={rectangle, rounded corners=2pt,
                draw=#1!70, fill=#1!9,
                text width=8cm, align=left,
                minimum height=0.72cm, inner sep=5pt,
                font=\footnotesize},
  num/.style={font=\footnotesize\bfseries, text=#1},
  node distance=0.22cm
]
\node[layer=CKernel]  (l1)  {\textbf{1. Identity}\enspace Agent DID verification and signature check};
\node[layer=CKernel,  below=of l1] (l2)  {\textbf{2. Capability}\enspace UCAN-style tool and memory access grants};
\node[layer=CExec,    below=of l2] (l3)  {\textbf{3. Budget}\enspace Token, memory, time, and tool-call limits};
\node[layer=CExec,    below=of l3] (l4)  {\textbf{4. Safety}\enspace Invariant checks, bounds, poison-pill detection};
\node[layer=CLogic,   below=of l4] (l5)  {\textbf{5. Governance}\enspace Pre/post conditions from CLS contracts};
\node[layer=CLogic,   below=of l5] (l6)  {\textbf{6. Compliance}\enspace Regulatory rules — HIPAA, GDPR, SOC2};
\node[layer=CProto,   below=of l6] (l7)  {\textbf{7. Audit}\enspace Logging requirements and retention policies};
\node[layer=CProto,   below=of l7] (l8)  {\textbf{8. Rate Limit}\enspace Per-agent and per-tenant throttling};
\node[layer=CDX,      below=of l8] (l9)  {\textbf{9. Quota}\enspace Tenant-level resource ceilings};
\node[layer=CDX,      below=of l9] (l10) {\textbf{10. Policy}\enspace Organisation-specific custom rules};
\node[layer=CDanger, below=of l10] (l11) {\textbf{11. Emergency}\enspace Circuit breaker and kill switch};

\node[right=0.3cm of l1, font=\tiny\itshape, text=CKernel] {always-on};
\node[right=0.3cm of l11, font=\tiny\itshape, text=CDanger] {last resort};

\end{tikzpicture}
\caption{11-layer predicate enforcement chain — evaluated in order for every kernel operation}
\label{fig:predicate-chain}
\end{figure}

\section{Cryptographic Receipt Chain}

Every session close produces a signed \textbf{ExecutionReceipt}. Receipts are
themselves stored as \texttt{Evidence} MemPackets, so the full history of any
agent's execution is a chain of content-addressed, signed receipts — each
referencing the CID of the previous one.

\begin{figure}[H]
\centering
\begin{tikzpicture}[
  receipt/.style={rectangle, rounded corners=4pt,
                  draw=CKernel!70, fill=CKernel!8,
                  text width=3cm, align=center,
                  minimum height=1.4cm, font=\footnotesize,
                  inner sep=5pt},
  node distance=1.4cm
]
\node[receipt] (r1) at (0,0)   {Receipt\\\texttt{cid: R-001}\\sig: Ed25519};
\node[receipt] (r2) at (3.6,0) {Receipt\\\texttt{cid: R-002}\\parent: R-001};
\node[receipt] (r3) at (7.2,0) {Receipt\\\texttt{cid: R-003}\\parent: R-002};
\node[rectangle, draw=gray!50, fill=gray!5, rounded corners=3pt,
      text width=2cm, align=center, minimum height=1.4cm, font=\footnotesize]
  (r4) at (10.2,0) {\ldots\\(chain head)};

\draw[arr=CKernel] (r1) -- (r2);
\draw[arr=CKernel] (r2) -- (r3);
\draw[arr=CKernel, dashed] (r3) -- (r4);

\node[below=0.2cm of r1, font=\footnotesize\itshape, text=gray]
  {session 1};
\node[below=0.2cm of r2, font=\footnotesize\itshape, text=gray]
  {session 2};
\node[below=0.2cm of r3, font=\footnotesize\itshape, text=gray]
  {session 3};

\end{tikzpicture}
\caption{Execution receipt chain — each receipt CID-references its predecessor, forming a tamper-evident audit log}
\label{fig:receipt-chain}
\end{figure}

To verify that an agent's history has not been tampered with, an auditor needs
only the chain head CID and the raw content of each receipt. No trusted third
party is required.

\section{Storage Backends}

The kernel abstracts over multiple storage backends through a single
\texttt{MemoryStore} trait. All backends provide the same read/write/query
semantics; only durability and scale characteristics differ:

\begin{table}[H]
\centering
\renewcommand{\arraystretch}{1.4}
\begin{tabularx}{\textwidth}{@{}llX@{}}
\toprule
\textbf{Backend} & \textbf{Use case} & \textbf{Characteristics} \\
\midrule
SQLite       & Local development, single-node & Zero-config, embedded, ACID \\
PostgreSQL   & Production, multi-node         & Full replication, WAL, high throughput \\
S3-compatible & Large binary content           & Object storage, unlimited scale \\
IPFS         & Decentralised deployments      & Content-addressed natively, p2p \\
\bottomrule
\end{tabularx}
\caption{Supported storage backends — all implement the same MemoryStore interface}
\end{table}
